% ============================================================
% Formulario, argomento 4: condizionamento e indipendenza (dispensa, capitoli 3 e 4)
% ============================================================
\argomento{Probabilità condizionata e indipendenza}{Dispensa: capitoli 3 e 4}

\voce{condizionata}{Probabilità condizionata}{probabilità di $A$ sapendo che si è verificato $B$}
  {P(A \mid B) = \frac{P(A \cap B)}{P(B)} \qquad (P(B) > 0)}
  {\variante{spazio equiprobabile}{P(A \mid B) = \frac{\card{A \cap B}}{\card{B}}}
   \variante{con le frequenze}{f_n(A \mid B) = \frac{n_{A \cap B}}{n_B}}
   \variante{è una probabilità: valgono tutte le regole, per esempio}{P(\compl{A} \mid B) = 1 - P(A \mid B)}}

\voce{composta}{Legge della probabilità composta}{probabilità di un'intersezione}
  {P(A \cap B) = P(B \mid A)\, P(A) = P(A \mid B)\, P(B)}
  {\variante{tre eventi \oltre}{P(A \cap B \cap C) = P(A)\, P(B \mid A)\, P(C \mid A \cap B)}}

\voce{partizione}{Partizione di $\Omega$}{i «casi»: ne succede uno e uno solo}
  {\bigcup_{i=1}^{M} E_i = \Omega \qquad E_i \cap E_j = \emptyset \ \ (i \neq j)}
  {\variante{la più semplice}{\{B, \compl{B}\}}
   \variante{conseguenza}{\sum_{i=1}^{M} P(E_i) = 1}}

\voce{totale}{Legge della probabilità totale}{$\{E_i\}$ partizione di $\Omega$}
  {P(A) = \sum_{i=1}^{M} P(A \cap E_i) = \sum_{i=1}^{M} P(A \mid E_i)\, P(E_i)}
  {\variante{due casi, $B$ e $\compl{B}$}{P(A) = P(A \mid B)\, P(B) + P(A \mid \compl{B})\, P(\compl{B})}}

\voce{bayes}{Teorema di Bayes}{$P(A) > 0$, $P(B) > 0$}
  {P(B \mid A) = P(A \mid B)\, \frac{P(B)}{P(A)}}
  {\variante{con una partizione $\{E_i\}$}
            {P(E_i \mid A) = \frac{P(A \mid E_i)\, P(E_i)}{\sum_{j=1}^{M} P(A \mid E_j)\, P(E_j)}}
   \variante{due casi, $B$ e $\compl{B}$}
            {P(B \mid A) = \frac{P(A \mid B)\, P(B)}{P(A \mid B)\, P(B) + P(A \mid \compl{B})\, P(\compl{B})}}}

\voce{indipendenza}{Eventi indipendenti}{}
  {P(A \cap B) = P(A)\, P(B)}
  {\variante{equivalente, se $P(B) > 0$}{P(A \mid B) = P(A)}
   \variante{anche i complementari sono indipendenti}{P(\compl{A} \cap \compl{B}) = P(\compl{A})\, P(\compl{B})}
   \variante{e anche $A$ e $\compl{B}$}{P(A \cap \compl{B}) = P(A)\, P(\compl{B})}
   \variante{$n$ prove indipendenti}{P(A_1 \cap A_2 \cap \cdots \cap A_n) = P(A_1)\, P(A_2) \cdots P(A_n)}}
